\section{Cálculo algebraico de equilibrios y espectros}
\label{app:algebra-espectros}

Para un campo polinomial con parámetros racionales, las ecuaciones
\(f(p)=0\), el jacobiano \(J_f(p)\) y el polinomio
\(p_p(\lambda)=\det(\lambda I-J_f(p))\) pueden obtenerse mediante
operaciones exactas. La representación racional de un decimal finito conserva
el parámetro declarado; el redondeo se introduce al aproximar las raíces.
Para evaluar una aproximación \(\widetilde p\) y sus valores propios
\(\widetilde\lambda_j\), se consideran
\[
 r_f=\|f(\widetilde p)\|,\qquad
 r_j=|p_{\widetilde p}(\widetilde\lambda_j)|,\qquad
 r_{\mathrm{tr}}=
 \left|\sum_j\widetilde\lambda_j-\operatorname{tr}J_f(\widetilde p)\right|.
\]
Un residuo pequeño indica consistencia numérica. La existencia y unicidad de
una raíz próxima requieren una cota adicional, por ejemplo mediante un
teorema de Newton validado; el residuo aislado no proporciona esa cota.

En el Lorenz generalizado con
\((\sigma,r,a)=(3.4,6.8,-0.5)\), los polinomios característicos son
\begin{align*}
 p_0(\lambda)&=\lambda^3+5.4\lambda^2-15.32\lambda-19.72,\\
 p_{\pm}(\lambda)&=\lambda^3+5.4\lambda^2
 +14.8476942706167\lambda+78.88.
\end{align*}
El coeficiente lineal de \(p_\pm\) se muestra redondeado. La suma espectral
es \(-5.4\), igual a la traza de los tres jacobianos. Para un sistema Caputo
conmensurable de orden \(q\), si ningún autovalor es cero, el margen angular
\[
 m_q(p)=\min_j|\arg\lambda_j|-\frac{q\pi}{2}
\]
determina la estabilidad asintótica del sistema lineal cuando es positivo.
Su aplicación al campo no lineal requiere las hipótesis de regularidad y
linealización establecidas en los fundamentos.

Las realizaciones de Lorenz generalizado y Rabinovich--Fabrikant se comprueban
mediante las identidades
\[
 AX+B\Phi(X)-f(X)=0,\qquad
 (zI-A)W(z)-B=0,
\]
en el dominio \(\det(zI-A)\ne0\). Estas igualdades justifican las
transferencias empleadas en el balance MIMO. La solución del balance y la
integración desde sus semillas constituyen problemas posteriores.
